\documentclass[10pt,a4paper]{article}
\usepackage[utf8]{inputenc}
\usepackage[T1]{fontenc}
\usepackage[ngerman]{babel}
\usepackage[margin=1.7cm,top=1.5cm,bottom=1.7cm]{geometry}
\usepackage{titlesec}
\titleformat*{\section}{\normalfont\large\bfseries}
\titleformat*{\subsection}{\normalfont\normalsize\bfseries}
\titlespacing*{\section}{0pt}{8pt}{3pt}
\titlespacing*{\subsection}{0pt}{6pt}{2pt}
\usepackage{amsmath,amssymb}
\usepackage{enumitem}
\usepackage{listings}
\usepackage{xcolor}
\usepackage{array}
\usepackage{fancyhdr}
\usepackage{parskip}
\usepackage{needspace}

% Ankreuzkästchen
\newcommand{\mcbox}{\raisebox{-1pt}{\fbox{\rule{0pt}{7pt}\rule{7pt}{0pt}}}\hspace{6pt}}

% Antwortlinien: n Linien über die volle Breite
\newcommand{\antwortlinien}[1]{%
  \par\vspace{2pt}%
  \newcount\zeilen \zeilen=#1
  \loop
    \noindent\rule{\linewidth}{0.4pt}\par\vspace{11pt}%
    \advance\zeilen by -1
  \ifnum\zeilen>0 \repeat
}

\lstset{
  language=Python,
  basicstyle=\ttfamily\footnotesize,
  keywordstyle=\bfseries,
  commentstyle=\itshape\color{gray},
  showstringspaces=false,
  frame=single,
  framerule=0.3pt,
  xleftmargin=8pt,
  framexleftmargin=6pt,
  aboveskip=6pt,
  belowskip=6pt
}

\pagestyle{fancy}
\fancyhf{}
\fancyhead[L]{\small MDT5/2 -- Session 1: Einführung \& Normalisierung}
\fancyhead[R]{\small Übungsaufgaben zur Klausurvorbereitung}
\fancyfoot[C]{\small Seite \thepage}
\renewcommand{\headrulewidth}{0.4pt}

\setlist[itemize]{itemsep=2pt,topsep=2pt}

\begin{document}

\begin{center}
  {\Large\bfseries Übungsaufgaben Session 1}\\[4pt]
  {\large Kapitel 1 (Einführung) \& Kapitel 2 (Normalisierung), Praktika 01a/01b}\\[4pt]
  \small Stil der Übungssammlung MDT5/2 -- generierte Aufgaben, keine Originale
\end{center}

\vspace{4pt}

\section*{Aufgabe 1 -- Multiple Choice mit Begründung}

\emph{Markieren Sie in jeder Teilaufgabe die einzige richtige von drei Aussagen und begründen Sie Ihre
Entscheidung. Punkte gibt es nur mit zutreffender Begründung.}

\subsection*{1.1 \ Für die Lernparadigmen gilt:}

\mcbox Beim Deep Learning werden Vorverarbeitung, Merkmalsextraktion und Klassifikation
vollständig aus Daten gelernt.

\mcbox Beim klassischen Maschinellen Lernen (kein Deep Learning) werden die Merkmale manuell
festgelegt und nur die Klassifikationsregeln aus Daten gelernt.

\mcbox Expertensysteme benötigen große Mengen an Trainingsdaten, um ihre Regeln zu konfigurieren.

Begründung:
\antwortlinien{2}

\subsection*{1.2 \ Für Novelty Detection und Outlier Detection gilt:}

\mcbox Bei der Outlier Detection ist auch im Training unbekannt, was in den Daten normal und
was Anomalie ist.

\mcbox Novelty Detection benötigt im Training ausreichend gelabelte Beispiele beider Klassen.

\mcbox Outlier Detection setzt voraus, dass die Trainingsdaten ausschließlich Normaldaten enthalten.

Begründung:
\antwortlinien{2}

\subsection*{1.3 \ Für Anomalien gilt:}

\mcbox Eine Kontextanomalie weicht immer auch isoliert betrachtet von der Beschaffenheit aller
anderen Punkte ab.

\mcbox Ungewöhnlich hohe Ausgaben an Weihnachten sind ein Beispiel dafür, dass ein Punkt erst
durch seinen Kontext zur Anomalie wird oder eben keine ist.

\mcbox Anomalien sind per Definition Messfehler und sollten immer aus den Daten entfernt werden.

Begründung:
\antwortlinien{2}

\subsection*{1.4 \ Für die Normalisierung gilt:}

\mcbox Nach einer Z-Transformation (Standardisierung) liegen alle Werte garantiert im
Bereich $[-1;1]$.

\mcbox Die Min-Max-Normalisierung verwendet Minimum und Maximum, die auf den Trainingsdaten
bestimmt wurden.

\mcbox Der \texttt{StandardScaler} muss auf den Testdaten erneut gefittet werden, damit auch
diese Mittelwert 0 haben.

Begründung:
\antwortlinien{2}

\section*{Aufgabe 2 -- Kurzantworten (Kapitel 1)}

\Needspace*{6cm}
a) Skizzieren Sie das Mengendiagramm von Künstlicher Intelligenz, Maschinellem Lernen,
Deep Learning und wissensbasierten Systemen.

\vspace{4cm}

b) Nennen Sie die drei Stufen der Mustererkennungspipeline und je zwei typische Operationen
pro Stufe.
\antwortlinien{4}

c) Geben Sie die Anomalie-Definition nach Chandola et al. sinngemäß wieder.
\antwortlinien{2}

d) Warum ist Anomalieerkennung \textbf{kein} gewöhnliches überwachtes Lernen?
Nennen Sie mindestens zwei Gründe aus der Vorlesung.
\antwortlinien{3}

e) Nennen Sie die zwei Hauptanwendungen der Outlier Detection aus der Vorlesung.
\antwortlinien{2}

f) Was ist mit dem \glqq unbekannten Unbekannten\grqq{} gemeint, und warum spricht es für
Novelty Detection?
\antwortlinien{3}

\section*{Aufgabe 3 -- Normalisierung rechnen (Kapitel 2)}

Die Trainingsdaten eines Merkmals sind: \quad $4,\ 4,\ 8,\ 8$

a) Berechnen Sie Mittelwert und Standardabweichung.
\antwortlinien{2}

b) Min-Max-Normalisierung auf $[0;1]$: Transformieren Sie die neuen Werte $7$ und $10$.
\antwortlinien{2}

c) Z-Transformation: Transformieren Sie dieselben Werte $7$ und $10$.
\antwortlinien{2}

d) Einer der Werte aus b) verlässt den Bereich $[0;1]$. Warum kann das passieren und warum
ist das gerade bei Anomalieerkennung eher zu erwarten?
\antwortlinien{3}

e) Wie heißen die beiden scikit-learn-Klassen für diese zwei Normalisierungen?
\antwortlinien{1}

\section*{Aufgabe 4 -- Code beurteilen}

\emph{Beurteilen Sie beide Ausschnitte auf Fehlerfreiheit bezüglich der Normalisierung und
begründen Sie jeweils Ihre Entscheidung.}

\Needspace*{5.5cm}
\begin{lstlisting}
# Variante A
scaler = StandardScaler()
scaler.fit(np.concatenate([train_x, test_x]))
train_n = scaler.transform(train_x)
test_n  = scaler.transform(test_x)
\end{lstlisting}
\antwortlinien{3}

\Needspace*{5.5cm}
\begin{lstlisting}
# Variante B
scaler = StandardScaler()
scaler.fit(train_x)
train_n = scaler.transform(train_x)
test_n  = scaler.transform(test_x)
\end{lstlisting}
\antwortlinien{3}

\Needspace*{13.5cm}
\section*{Aufgabe 5 -- Verfahrenslandkarte \ (= Start für Ihr handschriftliches Blatt \textcircled{\small 1})}

Ordnen Sie die folgenden Verfahren aus dem Gedächtnis in die Tabelle ein
(danach mit Folie 33 aus Kapitel 1 abgleichen). \textbf{Geklammerte Verfahren = Klammer-Regel:
einordnen können, aber NICHT klausurrelevant (nicht lernen):}

\begin{center}
k-NN, (LOF), Isolation Forest, (Matrix Profiles), Histogramm, Mahalanobis, KDE, GMM,\\
(PCA), (k-Means), Autoencoder, (VAE), f-AnoGAN, OC-SVM, SVDD,\\
Deep SVDD, GOAD, CutPaste
\end{center}

\vspace{6pt}
\renewcommand{\arraystretch}{1.2}
\begin{center}
\begin{tabular}{|>{\bfseries}p{1.6cm}|p{2.9cm}|p{2.9cm}|p{2.9cm}|p{2.9cm}|}
\hline
 & \textbf{Distanz} & \textbf{Probabilistisch} & \textbf{Rekonstruktion} & \textbf{Klassifikation} \\
\hline
Deep & & & & \\[2.2cm]
\hline
Shallow & & & & \\[3.2cm]
\hline
\end{tabular}
\end{center}

\section*{Aufgabe 6 -- NumPy-Verständnis (zu Praktikum 01b, ohne Rechner lösen!)}

Gegeben: \ \texttt{a = np.array([[1, 2, 3], [4, 5, 6]])}

a) Was ist \texttt{a.shape}?
\antwortlinien{1}

b) Was liefert \texttt{np.mean(a, axis=0)} -- Werte und Shape? Und \texttt{np.mean(a, axis=1)}?
\antwortlinien{2}

c) Was liefert \texttt{a + np.array([10, 20, 30])}? Wie heißt dieser Mechanismus?
\antwortlinien{2}

d) Was liefert \texttt{a[:, 1]}? Und \texttt{a[0, :]}?
\antwortlinien{1}

e) Was ist \texttt{a.reshape(3, 2)} -- und warum ist es \textbf{nicht} dasselbe wie \texttt{a.T}?
\antwortlinien{2}

\Needspace*{10cm}
\section*{Aufgabe 7 -- Transfer: Novelty oder Outlier Detection?}

\emph{Entscheiden Sie für beide Szenarien begründet, ob Sie das Problem als Novelty Detection
oder Outlier Detection aufsetzen würden. Stichwort: Datenlage beim Bau des Systems.}

a) Ein Herzspezialist gibt Ihnen 10\,000 EKG-Aufnahmen gesunder Patienten. Das System soll
später krankhafte Veränderungen melden -- auch Krankheitsbilder, die heute noch unbekannt sind.
\antwortlinien{4}

b) Sie erhalten einen großen, ungelabelten Datensatz mit Kreditkartentransaktionen und sollen
\glqq Auffälligkeiten finden, die sich lohnen könnten anzuschauen\grqq{} -- niemand kann Ihnen
sagen, was normal ist.
\antwortlinien{4}

\vspace{10pt}
\begin{center}
\footnotesize\itshape
Nach Bearbeitung: Antworten an Claude schicken (Foto genügt) -- Korrektur mit Begründungs-Check.
\end{center}

\end{document}
